\chapter{Introduction}
\label{chap:introduction}

\section{Background and Context}
\label{sec:intro_background}
Enterprise Information Management Systems (EIMS) in modern commercial hospitality and culinary operations operate at the critical nexus of real-time resource scheduling, transactional financial accounting, relational inventory governance, and multi-tenant organizational management \parencite{fowler2002patterns, chong2006multi}. In contemporary food service enterprises---ranging from independent dining establishments to multi-branch franchise networks and academic campus dining commissaries---information management platforms must reliably reconcile high-concurrency reservation intake with strict physical floor plan constraints, table capacities, kitchen preparation limits, and customer service workflows.

Historically, food and beverage operations relied primarily on manual physical reservation logbooks or isolated, single-tenant desktop databases. While manual logbooks are notoriously vulnerable to transcription errors, catastrophic double-bookings, untracked cancellations, and unmitigated revenue loss from customer no-shows, modern commercial Software-as-a-Service (SaaS) suites introduce distinct architectural drawbacks. Mainstream commercial hospitality software is typically delivered as monolithic, multi-gigabyte cloud subscriptions burdened by deep dependencies on third-party backend frameworks, complex client-side JavaScript runtimes, remote Content Delivery Networks (CDNs), and opaque telemetry services \parencite{stonebraker2007end, bezemer2010multitenant}. In localized, self-hosted, edge-deployed, or air-gapped evaluation environments---such as the rigorous computing testbeds mandated by the Department of Information Technology at Cavite State University (CvSU)---such external dependencies compromise architectural determinism, auditability, operational portability, and information privacy.

Furthermore, contemporary reservation systems routinely exhibit severe concurrency anomalies during peak booking spikes. When prospective diners or front-of-house staff concurrently attempt to reserve dining tables, naive relational designs that rely on unconstrained continuous temporal ranges or optimistic client-side updates suffer from race conditions, resulting in overlapping reservations, zone overcapacity, and customer dissatisfaction \parencite{bernstein1987concurrency, kung1981optimistic}. When large parties require combining adjacent physical tables, transactional complexity compounds significantly. Compounding these scheduling challenges, standard commercial systems frequently employ imprecise IEEE-754 floating-point numeric types for financial accounting, introducing cumulative rounding errors into customer deposit ledgers and loyalty point balances.

To address these systemic vulnerabilities, the Enterprise Information Management System (EIMS) was conceptualized, designed, and engineered as a clean-room, multi-tenant restaurant ERP and high-concurrency reservation platform for the DCIT 24 Information Management capstone at Cavite State University - Main Campus. Built strictly from first principles using native PHP 8.3 without external backend or frontend frameworks, EIMS demonstrates that industrial-grade multi-tenancy, provably correct reservation concurrency control, dual-driver relational resilience, and active data integrity guards can be achieved with total architectural self-sufficiency.

\section{Enterprise Restaurant Information Management Challenges}
\label{sec:intro_challenges}
Commercial restaurant and hospitality operations face multi-faceted data management challenges that span several distinct operational tiers:

\begin{enumerate}[label=\textbf{C\arabic*.}]
    \item \textbf{Peak-Hour Booking Concurrency and Physical Resource Contention:} Unlike static e-commerce inventories, dining tables represent finite, temporally bound physical assets. During high-demand reservation intervals (e.g., weekend dinner rushes or holiday seatings), hundreds of prospective diners simultaneously query availability and submit booking requests. Without strict transactional locking and boundary discretization, concurrent read-modify-write sequences generate severe double-booking anomalies \parencite{papadimitriou1979serializability, gray1993transaction}.
    \item \textbf{Dynamic Floor Plan Topologies and Combinable Seating:} Restaurant dining rooms are organized into distinct physical zones (e.g., Main Dining, Al Fresco, Private Function Rooms, Mezzanine) with varying acoustic, ambient, and capacity characteristics. Accommodating variable party sizes requires dynamic table combinations, where two or more adjacent physical tables are merged into a single logical seating unit. Scheduling combinable tables demands atomic multi-resource allocation: either all required tables are locked and reserved simultaneously, or the entire transaction is rolled back cleanly.
    \item \textbf{Kitchen Inventory Depletion and Menu Availability Synchronization:} A fundamental operational breakdown in hospitality management occurs when front-of-house staff accept orders for dishes whose raw ingredients have been exhausted in the kitchen. Maintaining relational synchronization between kitchen inventory stock levels, Bill of Materials (BOM) recipe requirements, and the real-time menu availability (the culinary ``86-list'') requires tight relational coupling across inventory movements and menu display states.
    \item \textbf{Dual-Ledger Financial Precision and Auditability:} Managing reservation commitment deposits and customer loyalty rewards requires absolute numerical precision. Floating-point arithmetic introduces microscopic rounding errors that accumulate across thousands of customer transactions, creating regulatory compliance liabilities and ledger discrepancies. A robust hospitality information system must implement immutable, append-only ledger accounting using exact integer arithmetic.
    \item \textbf{Multi-Tenant Governance and Data Privacy Compliance:} Modern enterprise platforms must support multiple independent restaurant entities within a single operational deployment. Ensuring complete cryptographic and logical data isolation between competing tenant organizations is paramount to prevent Insecure Direct Object References (IDOR), unauthorized cross-tenant data harvesting, and non-compliance with statutory data privacy mandates such as Republic Act No.~10173 (Data Privacy Act of 2012) \parencite{ra10173, owasp2021top10}.
\end{enumerate}

\section{Problem Statement}
\label{sec:intro_problem}
In evaluating existing hospitality management and restaurant booking software, four persistent engineering flaws compromise operational integrity, data safety, and software longevity:

\begin{enumerate}[label=\textbf{P\arabic*.}]
    \item \textbf{Vulnerability to Concurrency Race Conditions:} Conventional database schemas represent reservations with continuous start and end timestamps (\texttt{DATETIME}). Under standard transaction isolation levels without mathematical slot discretization, concurrent requests targeting the same table evaluate availability simultaneously, pass procedural overlap checks, and commit conflicting records, creating physical double-bookings.
    \item \textbf{Single-Point Persistence Failure in Constrained Environments:} Traditional web applications couple their data persistence exclusively to a centralized database daemon (e.g., MySQL). In local demonstration environments, air-gapped evaluation labs, or container chaos scenarios, any disruption in database daemon connectivity results in immediate, total application failure. The absence of an autonomous, zero-configuration local fallback persistence engine cripples operational resilience.
    \item \textbf{Relational Corruption via Unguarded Active Mutations:} When restaurant managers modify, downsize, or archive dining tables or floor plans while active, confirmed future reservations remain linked to those resources, systems create orphaned foreign keys or reduce table capacities below party sizes, breaking operational continuity between floor layout and customer commitments.
    \item \textbf{Software Bloat and Security Surface Expansion:} Modern enterprise stacks routinely incorporate hundreds of third-party dependencies (via Composer or npm), creating deep supply chain security vectors, introducing inline stylesheet violations that break Content Security Policy (CSP) directives, and enabling account enumeration attacks through disparate authentication timings.
\end{enumerate}

\section{Business Rationale and Operational Value}
\label{sec:intro_business_rationale}
The implementation of EIMS delivers measurable business, operational, and pedagogical value:

\begin{itemize}
    \item \textbf{Mitigation of Revenue Loss from Scheduling Failures:} By mathematically eliminating double-booking collisions and automating waitlist queue promotions when cancellations occur, EIMS maximizes dining table utilization rates and prevents the direct financial and reputational damages associated with turned-away patrons.
    \item \textbf{Air-Gapped Operational Self-Sufficiency:} By discarding external cloud APIs, CDN scripts, and third-party vendor frameworks, EIMS operates with 100\% local autonomy. A restaurant branch can execute full ERP and reservation scheduling operations on low-cost local hardware or an internal local area network (LAN) during external internet outages.
    \item \textbf{Complete Financial Auditability:} Signed 64-bit integer centavo ledgers guarantee exact monetary balances for booking deposit holds, forfeitures, and loyalty point redemptions, ensuring complete compliance with accounting standards and auditing requirements.
    \item \textbf{Pedagogical Excellence at Cavite State University:} Developed as the capstone deliverable for DCIT 24 Information Management, EIMS demonstrates the practical power of pure relational database theory (Third Normal Form normalization, composite foreign keys, transactional locking, and indexing strategies) in solving complex enterprise problems without relying on opaque framework abstractions.
\end{itemize}

\section{Document Structure}
\label{sec:intro_structure}
This technical project report adheres rigorously to the 10-section project documentation standard mandated by the Cavite State University DCIT 24 Information Management syllabus:

\begin{itemize}
    \item \textbf{Section 1: Title Page} establishes formal institutional, departmental, authorial, and academic identification.
    \item \textbf{Section 2: Introduction (Chapter~\ref{chap:introduction})} outlines the background, industry challenges, formal problem statement, business rationale, and structural overview.
    \item \textbf{Section 3: Project Objectives (Chapter~\ref{chap:project_objectives})} defines the general research objective alongside specific technical, architectural, and business milestones.
    \item \textbf{Section 4: Scope and Limitations (Chapter~\ref{chap:scope_limitations})} details the functional boundaries, supported user roles (SuperAdmin, Admin, FOH, Kitchen, Customer), implemented features, and explicit project delimitations.
    \item \textbf{Section 5: System Requirements (Chapter~\ref{chap:system_requirements})} specifies minimum and recommended hardware specifications, software runtimes, web server rewrite rules, and deployment network topologies.
    \item \textbf{Section 6: Entity-Relationship Diagram (Chapter~\ref{chap:erd})} formalizes the 18-table relational architecture across five functional domains, detailing Crow's foot cardinalities and composite keys.
    \item \textbf{Section 7: Database Design (Chapter~\ref{chap:database_design})} provides the exhaustive Data Dictionary for all database entities, mathematical normalization proofs (1NF through BCNF), indexing schemes, and dual-driver SQL dialect branching.
    \item \textbf{Section 8: System Features (Chapter~\ref{chap:system_features})} presents a comprehensive walkthrough of implemented system capabilities, including multi-tenant isolation, bitmask RBAC, the 12x8 floor plan grid, the \texttt{TableAllocator} engine, waitlist promotions, kitchen 86-lists, and dual-ledger accounting.
    \item \textbf{Section 9: Testing and Results (Chapter~\ref{chap:testing_results})} details the empirical verification methodology, reporting outcomes from 41 integration tests, 82 security hardening tests, 15 adversarial Red Team probes, the formal double-booking prevention proof, stress benchmarks, and chaos failover tests.
    \item \textbf{Section 10: Conclusion (Chapter~\ref{chap:conclusion})} synthesizes project achievements against defined objectives, analyzes operational impact, verifies curriculum alignment, and proposes future engineering enhancements.
    \item \textbf{Individual Member Reflections (Chapter~\ref{chap:reflections})} documents the reflective practice, individual accountability, technical contributions, database concepts mastered, and personal evaluations of all eight project team members.
\end{itemize}
